\documentclass[11pt]{article}
\usepackage{amsmath,amssymb,amsthm,hyperref}
\newtheorem{theorem}{Theorem}
\newtheorem{lemma}{Lemma}
\begin{document}
\title{An exact-value upper bound from a reciprocal Segre degeneration}
\author{}
\date{}
\maketitle

\paragraph{Status.}
This is a complete written deduction using the external secant theorem
of Catalisano, Geramita and Gimigliano stated below. It is not a Lean
proof: that external theorem and the analytic rank argument have not
been imported into the formal development. The consequence improves
the constant in the existing order $2^n/n$ upper bound. It gives
neither a logarithmic saving nor the proposed matching lower bound.
The original input dimension, atom admissibility, $H^*$ and
\texttt{worstCaseHeadComplexity} are unchanged.

\begin{theorem}\label{thm:exact}
For every integer $n\geq1$, put $N=2^n$ and
$R=\lceil N/(n+1)\rceil$. Every function
$y:\{0,1\}^n\to\mathbb R$ is an exact sum of $2R$ fractions
\[
 y(x)=\sum_{h=1}^{2R}\frac{c_h}{1+\sum_{i=1}^n b_{hi}x_i},
 \qquad b_{hi}>0.
\]
Consequently, in the unchanged HeadComplexity model,
\[
 H^*(f)\leq2\left\lceil\frac{2^n}{n+1}\right\rceil,
 \qquad
 \operatorname{worstCaseHeadComplexity}(n)
 \leq2\left\lceil\frac{2^n}{n+1}\right\rceil.
\]
\end{theorem}

\paragraph{External theorem used.}
Let $X_n$ be the Segre image of $(\mathbb P^1_{\mathbb C})^n$ in
$\mathbb P^{2^n-1}_{\mathbb C}$. Theorem 4.1 of Catalisano, Geramita
and Gimigliano says that its $s$th secant variety has dimension
\[
 \min\{s(n+1)-1,2^n-1\},
\]
except that the third secant of $X_4$ has dimension $13$.
In particular the $R$th secant fills the ambient projective space.
The exception does not affect $R$: when $n=4$, $R=4$.
See \emph{J. Algebraic Geometry} 20 (2011), 295--327,
\href{https://doi.org/10.1090/S1056-3911-10-00537-0}{doi:10.1090/S1056-3911-10-00537-0},
or the authors' \href{https://arxiv.org/abs/0809.1701}{arXiv:0809.1701}.
The only consequence needed is that the polynomial map
\[
 \Psi(c,v)_S=\sum_{h=1}^R c_h\prod_{i\in S}v_{hi},
 \qquad S\subseteq[n],
\]
has a Jacobian minor of size $N$ that is not the zero polynomial.
Indeed, the affine chart $\prod_i(1,v_i)$ is Zariski dense in the
Segre variety, so the secant theorem makes $\Psi$ dominant.
In characteristic zero, dominance of this polynomial map implies
that its differential has rank $N$ on a nonempty Zariski open set.
Equivalently, this follows from the tangent-space statement of
Terracini's lemma, Proposition 2.2 in the same paper.

\begin{lemma}[Positive real parameters]\label{lem:positive}
There are real parameters $(c^*,v^*)$ with every $c^*_h>0$ and
every $v^*_{hi}>0$ at which $D\Psi$ has rank $N$.
\end{lemma}
\begin{proof}
The nonzero Jacobian minor supplied above is a polynomial with real
(in fact integer) coefficients. A nonzero real polynomial cannot
vanish on a nonempty real open box. Apply this to any open box
contained in the strictly positive orthant of all parameters.
Complex generic rank is being used only to establish that this
real polynomial is nonzero; no assertion about arbitrary real
tensor decompositions is needed.
\end{proof}

\begin{lemma}[Differentiable degeneration]\label{lem:degeneration}
For $v\in(0,\infty)^n$ and $\epsilon>0$, write
\[
 K_\epsilon(v)(x)=\frac1{1+\epsilon\sum_i v_i x_i}.
\]
There is an invertible real linear map
$T_\epsilon:\mathbb R^{\{0,1\}^n}\to
\mathbb R^{\mathcal P([n])}$ such that
\[
 T_\epsilon K_\epsilon(v)\longrightarrow
       \left(\prod_{i\in S}v_i\right)_{S\subseteq[n]}
\]
in $C^1$ on every compact subset of $(0,\infty)^n$ as
$\epsilon\downarrow0$.
\end{lemma}
\begin{proof}
For a table $F$, let its Boolean Moebius coefficient be
\[
 (\mathcal M F)_S
   =\sum_{T\subseteq S}(-1)^{|S|-|T|}F(\boldsymbol1_T).
\]
The inverse is $F(\boldsymbol1_T)=\sum_{S\subseteq T}(\mathcal M F)_S$.
Define
\[
 (T_\epsilon F)_S
   =\frac{(\mathcal M F)_S}{(-1)^{|S|}|S|!\epsilon^{|S|}}.
\]
Thus $T_\epsilon$ is invertible for each fixed $\epsilon>0$.
For $k=|S|$, iterated one-variable fundamental theorems of calculus
give the exact identity
\[
 (T_\epsilon K_\epsilon(v))_S
 =\left(\prod_{i\in S}v_i\right)
   \int_{[0,1]^S}
      \left(1+\epsilon\sum_{i\in S}v_i t_i\right)^{-k-1}\,dt.
\]
For $S=\varnothing$ the product and integral are both $1$.
For nonempty $S$, the identity follows by integrating the mixed
derivative
\[
 \partial_S\frac1{1+\epsilon\sum_i v_i x_i}
 =(-1)^k k!\epsilon^k
     \frac{\prod_{i\in S}v_i}
          {(1+\epsilon\sum_i v_i x_i)^{k+1}}
\]
over the indicated coordinate face. On a compact positive parameter
set, the integral converges uniformly to $1$, and each of its first
$v$-derivatives converges uniformly to $0$. This follows directly
by differentiating the displayed integrand: the differentiated
integrand has an explicit factor $\epsilon$ and its denominator
is at least $1$. There are only finitely many $S$, proving the
claimed $C^1$ convergence.
\end{proof}

\begin{proof}[Proof of Theorem~\ref{thm:exact}]
Consider the smooth map on the open positive parameter orthant
\[
 \Phi_\epsilon(c,v)=\sum_{h=1}^R c_h K_\epsilon(v_h)
       \in\mathbb R^N.
\]
By Lemma~\ref{lem:degeneration}, $T_\epsilon\Phi_\epsilon$
converges in $C^1$ near $(c^*,v^*)$ to $\Psi$. Choose a Jacobian
minor that is nonzero at $(c^*,v^*)$. Its determinant remains
nonzero for all sufficiently small positive $\epsilon$.
Fix one such $\epsilon$. Since $T_\epsilon$ is an invertible
linear map, $D\Phi_\epsilon(c^*,v^*)$ itself has rank $N$.

For completeness, this step requires no norm bound on
$T_\epsilon$ as $\epsilon\downarrow0$. The limit is used only
to locate one nonzero determinant at one strictly positive,
fixed $\epsilon$; multiplication by that fixed invertible map
then preserves rank exactly.

The inverse function theorem, applied after fixing all parameters
outside the chosen $N$ columns of the differential, shows that
the image of $\Phi_\epsilon$ contains an open neighborhood of
$F_0=\Phi_\epsilon(c^*,v^*)$. The parameter neighborhood can
be chosen inside the positive orthant. For any prescribed
$y\in\mathbb R^N$, choose a positive scalar $\delta$ sufficiently
small that $F_0+\delta y$ belongs to this image. Hence there are
positive $c'_h,v'_{hi}$ such that
\[
 y(x)=\frac1\delta
  \left(\sum_{h=1}^R
      \frac{c'_h}{1+\epsilon\sum_i v'_{hi}x_i}
      -\sum_{h=1}^R
      \frac{c^*_h}{1+\epsilon\sum_i v^*_{hi}x_i}\right).
\]
This is exactly the required $2R$-head representation. There is
no limiting approximation in this final equality, no dimension
padding, and no restriction on the size of the final numerators.

Each denominator has constant coefficient $1$ and strictly positive
slopes, so it is positive on the whole Boolean cube. Each numerator
is constant, hence affine. The existing model conversion
\texttt{FracAtom.ofPositiveAffineRatio} realizes each such ratio
by one original atom; negative numerator constants are allowed.
Apply the representation to $y(x)=1$ when $f(x)$ is true and
$y(x)=-1$ when it is false. The signs are strict, the scalar bias
is zero, and
\texttt{computableWithHeadsN\_of\_positiveAffineRatios}
gives the claimed $H^*$ bound. Taking the finite maximum over
Boolean functions gives the bound against the unchanged
\texttt{worstCaseHeadComplexity}.
\end{proof}

\paragraph{Checks and limitations.}
The companion Python check computes ranks over a prime field at
explicit positive rational reciprocal parameters for $1\leq n\leq6$.
A full rank in this reduction certifies that the corresponding
real rational Jacobian has full rank. These finite checks test the
formulas; the argument for all $n$ uses the external secant theorem.
There is no claim that a fixed $\epsilon$ works uniformly in $n$,
or that the parameter sizes or the open image radius are controlled.
Neither is required in the unchanged model. A logarithmic improvement
would require a new strict-sign argument below this exact-value
rank scale; the degeneration and the external theorem do not provide
such an argument.
\end{document}
